\originalchapter{22.02 作业六:两核子体系}
题源: Spring 2012 Problem Set 6. 本章全部为编者独立解答. 原题日期存在排印问题; 物理参数按 PDF 原题保留.
\begin{exercise}\label{ex:2202-6-1}
原题 1. (a) 用 $\mathbf R=(m_1\mathbf x_1+m_2\mathbf x_2)/M$, $\mathbf r=\mathbf x_1-\mathbf x_2$, $\mathbf P=\mathbf p_1+\mathbf p_2$, $\mathbf p=(m_2\mathbf p_1-m_1\mathbf p_2)/M$ 证明 $E_{\rm kin}=P^2/(2M)+p^2/(2\mu)$. (b) 对 $H=p_1^2/(2m_1)+p_2^2/(2m_2)+V(|\mathbf x_1-\mathbf x_2|)$ 分离为两个微分方程, 质心动能是多少?
\end{exercise}
\begin{solution}
反解两动量后平方, 混合项系数为 $1/M-1/M=0$, 得所求. 取 $\Psi(\mathbf R,\mathbf r)=\e^{\ii\mathbf P\cdot\mathbf R/\hbar}\phi(\mathbf r)$, 第一方程 $-\hbar^2\nabla_R^2\Psi_R/(2M)=E_{\rm CM}\Psi_R$, 第二方程 $[-\hbar^2\nabla_r^2/(2\mu)+V]\phi=E_{\rm rel}\phi$. $E_{\rm CM}=P^2/(2M)$ 为任意非负连续值, 若选质心静止才为零.
\end{solution}
\begin{exercise}\label{ex:2202-6-2}
原题 2. 以能级相对井深的估算说明氘核没有束缚激发态, 分 (a) 更高径向态, (b) 非零 l; (c) 如果核力范围 $R_*>R_0$, 怎样使 $l=1$ 激发态存在而无高阶径向态, 给范围上下限.
\end{exercise}
\begin{solution}
采用同题组球方阱 $V_0=35$ MeV、$\mu c^2\simeq469.5$ MeV、$R_0=2.1$ fm. 最大井内波数 $k_0=\sqrt{2\mu V_0}/\hbar\simeq0.919\,\mathrm{fm}^{-1}$, 所以 $k_0R_0\simeq1.93$. 第二 $s$ 态阈值为 $3\pi/2=4.71$, 首个 $p$ 态为 $\pi=3.14$, 均未达到. 对 $p$ 阈值, 零能外径向 $u\propto r^{-1}$, 以内 $u\propto rj_1(k_0r)$ 匹配可化为 $j_0(k_0R)=0$, 第一根 $\pi$. 因而
$\pi/k_0<R_*<3\pi/(2k_0)$, 数值 $3.42<R_*<5.13$ fm. 端点恰为零能临界, 不计作正规束缚态. 这是指定模型的结论, 不用一个中心方阱证明真实核力所有通道的谱.
\end{solution}
\begin{exercise}\label{ex:2202-6-3}
原题 3. $u=A\sin kr+B\cos kr$ ($r<R_d$), $u=C\e^{-\kappa r}+D\e^{\kappa r}$ ($r>R_d$). (a) 用连续性、归一化求系数, 给 $V_0=35$ MeV、$R_d=2.1$ fm、E=$-2.2$ MeV. (b) 求均方根半径. (c) 求 $r>R_d$ 概率.
\end{exercise}
\begin{solution}
原点正规与远处可积给 $B=D=0$, 连续性给 $C=A\sin(kR_d)\e^{\kappa R_d}$, A 由式 \eqref{eq:deuteron-norm}. 
\begin{equation*}
 \langle r^2\rangle=A^2\int_0^{R_d}r^2\sin^2kr\dd r+C^2\int_{R_d}^\infty r^2\e^{-2\kappa r}\dd r.
\end{equation*}
外部部分, 外部积分为 $\e^{-2\kappa R_d}[R_d^2/(2\kappa)+R_d/(2\kappa^2)+1/(4\kappa^3)]$.

给定三参数不是精确相容的本征态: 用 $\mu c^2=469.5$ MeV 与 $\hbar c=197.327$ MeV fm, $k\cot(kR_d)+\kappa\simeq-0.0417\,\mathrm{fm}^{-1}$. 若坚持 E=$-2.2$ MeV, 可按上述连续且归一的近似波形计算, 但不能声称导数连续; 此版本 $A=0.5597$ fm$^{-1/2}$, $P_{\rm out}=0.6218$, 相对坐标 RMS $=3.8876$ fm. 固定 $V_0,R_d$ 并数值解匹配条件, 得相容 $B_d\simeq2.714$ MeV, $k=0.8824$, $\kappa=0.2558$ fm$^{-1}$, $A\simeq0.5769$ fm$^{-1/2}$, $P_{\rm out}\simeq0.6001$, 核子间距 RMS $\simeq3.612$ fm. 若指每核子相对于质心的 RMS 半径, 等质量近似为一半约 $1.806$ fm. 两套参数口径不能混用.
\end{solution}
\begin{exercise}\label{ex:2202-6-4}
原题 4. 说明双中子不束缚: (a) 哪个 l 能量最低可能束缚? (b) 对应轨道态交换性质? (c) 自旋态必须怎样, 结合核力自旋依赖说明总势能.
\end{exercise}
\begin{solution}
低能是 $l=0$, 没有离心抬升; 此空间态在交换时对称. 两中子的同位旋部分也对称, 费米总态反对称迫使自旋单重态. 氘核的主导三重态吸引不能直接适用, 单重态实验相互作用不足以形成自由双中子束缚核. Pauli 只限定通道, ``该通道不足以束缚''是经验核力输入; 若任意增大该通道吸引, 数学上完全可以束缚, 所以不能只由统计原理推出一切可能相互作用下都不存在.
\end{solution}
